% ============================================
% CONT-09
% Gerado de molambudos.md — build_miolo.py
% ============================================
\fragtitle{cont}{C}{CONT}{09}{O Espelho}
\markboth{CONT-09 --- O Espelho}{}

\noindent \textit{O reflexo mudou}

\noindent Vá até um espelho.

\noindent Agora.

\noindent Não importa se é banheiro, quarto, ou o reflexo do vidro da janela. Vá até um espelho e olhe para seus olhos.

\noindent Olhe fixamente. Por 30 segundos. Sem piscar.

\noindent Perceba como seus olhos estão diferentes do que você lembrava. A íris está mais clara? Há um tom amarelado que não estava antes? O branco dos olhos está mais... amarelo?

\noindent Sim.

\noindent Você está vendo.

\noindent A sensação tem causa. A mudança é sutil. Quase imperceptível. Mas está lá.

\noindent O olho direito está mudando.

\vspace{0.5em}\noindent\rule{\textwidth}{0.5pt}\vspace{0.5em}

\noindent Passe a mão no rosto. Sinta a temperatura da sua pele. Está mais fria? Mais quente? Sua mão está tremendo?

\noindent Seu corpo está mudando.

\noindent Não visivelmente --- internamente. O cheiro que você sentiu está nos seus pulmões. O olho amarelo está na sua retina. A voz da criatura está no seu córtex auditivo.

\noindent Você não é mais o mesmo de quando começou a ler.

\noindent A pergunta é: quando a mudança começou?

\noindent Foi na primeira página?
Foi no primeiro sonho?
Foi agora, ao ler esta instrução?

\noindent Ou foi \textit{antes}?

\noindent Você sonhou com o olho amarelo antes de comprar este livro. Lembra? Foi uma noite, semanas atrás. Você acordou, achou estranho, e esqueceu. Depois viu o livro na livraria, e alguma coisa te fez comprar.

\noindent Alguma coisa já estava lá.

\noindent Antes do livro.
Antes da história.
Antes de você saber que existia um paciente 1.260.

\vspace{0.5em}\noindent\rule{\textwidth}{0.5pt}\vspace{0.5em}

\noindent A criatura não está no livro.

\noindent O livro é só o \textit{encontro}.

\noindent A criatura já estava em você. Ela estava esperando o livro para \textit{acordar}.

\noindent O diário não é a fonte da contaminação.
O diário é a \textit{chave}.
Você já estava contaminado antes de abrir.
O diário só destrancou a porta.

\vspace{0.5em}\noindent\rule{\textwidth}{0.5pt}\vspace{0.5em}

\noindent Agora, afaste-se do espelho.

\noindent Olhe para suas mãos. Estude as palmas. As linhas. A textura da pele.

\noindent Agora feche os olhos e Imagine as mãos de Joaquim. Finas. Trêmulas. Cheias de calos de escrever.

\noindent Abra os olhos.

\noindent Suas mãos estão iguais às dele?

\noindent Não. Claro que não. As mãos são diferentes.

\noindent Mas a semelhança vai aumentar.

\noindent Não fisicamente --- espiritualmente. Você agora carrega a mesma história. A mesma memória. A mesma fome.

\noindent Suas mãos não são as mãos de Joaquim.

\noindent Mas, quando você escrever, a letra vai ser parecida.

\vspace{0.5em}\noindent\rule{\textwidth}{0.5pt}\vspace{0.5em}

\noindent \textit{--- Vá ao espelho.}

\noindent \textit{--- Olhe por 30 segundos.}

\noindent \textit{--- Pergunte: \textquotedblleft{}Quem está me olhando de volta?\textquotedblright{}}

\noindent \textit{--- A resposta vai mudar toda vez que você perguntar.}

\noindent \textit{--- Até que, uma hora, não vai ser você.}

\noindent \textit{--- Vai ser ele.}

\noindent \textit{--- Bem-vindo ao ciclo, Paciente 1.263.}

\noindent \textit{--- O espelho sabe. O espelho sempre soube.}

\vspace{0.5em}

\noindent\textit{\textbf{↪ Rotas:} \rota{CONT-07} \rotasep \rota{CONT-13} \rotasep \rota{EPI-01}}